\Needspace{20\baselineskip}
\section{Marco Teórico}

\subsection{Arquitectura de datos moderna y capas semánticas}

La implementación de sistemas de Inteligencia de Negocios Generativa (GenBI)
necesita separar la información de las operaciones diarias de aquella usada
para el análisis. Las bases de datos transaccionales, diseñadas para el
registro rápido y seguro de datos crudos, no contienen la lógica del negocio.
Conectar modelos de lenguaje directamente a estas bases de datos aumenta el
riesgo de generar métricas incorrectas o respuestas inventadas, ya que la
Inteligencia Artificial desconoce las reglas propias de la organización
\citep{negash2004fundamentos}. Investigaciones recientes señalan que la
complejidad de las tablas, los nombres poco claros de las columnas y la falta
de contexto son las principales causas de error en los sistemas actuales que
traducen texto a SQL \citep{rissaki2024agentic}.

Para solucionar este problema, la arquitectura de datos moderna utiliza
almacenes analíticos y construye una capa semántica, que funciona como un
puente de comunicación. Esta capa se encarga de traducir la complejidad
técnica de las tablas y sus relaciones en conceptos de negocio claros y
estandarizados. Estudios actuales sugieren automatizar la creación de esta
capa utilizando sistemas multi-agente. En este esquema, los modelos de
lenguaje organizan la información relacional en estructuras mucho más
fáciles de comprender, logrando una alta precisión al integrar los datos
de la empresa \citep{guerbas2025multiagent}.

Dentro del análisis de datos, la capa semántica proporciona a la Inteligencia
Artificial un entorno de trabajo seguro y controlado. En lugar de requerir que
el modelo deduzca cómo relacionar múltiples tablas en bases de datos complejas,
el agente interactúa únicamente con información validada y simplificada
\citep{rissaki2024agentic}. Recientemente, este enfoque se ha integrado en
sistemas que emplean modelos de lenguaje para procesar consultas. En estos
casos, la capa semántica permite interpretar con mayor claridad la intención
del usuario antes de generar el código. Esto asegura la obtención de
resultados confiables y estrictamente alineados con las reglas de la
organización \citep{tursio2026deepcontext}.

\subsection{Conversión de texto a lenguaje SQL}

La conversión de texto a lenguaje estructurado (Text-to-SQL) es una herramienta
clave para facilitar el acceso a la información en las empresas. Permite que
usuarios sin conocimientos de programación puedan consultar bases de datos
relacionales utilizando lenguaje natural. Este proceso analiza la estructura y
el significado de las preguntas para entender qué datos se están solicitando
\citep{luo2024nl2sql}. Sin embargo, el principal desafío técnico es resolver la
ambigüedad de las palabras y entender exactamente qué quiere el usuario,
especialmente cuando las bases de datos tienen nombres o estructuras poco
descriptivas \citep{srivastava2026comparative}.

Para mejorar su funcionamiento, los sistemas actuales utilizan técnicas de
vinculación de esquemas (\textit{Schema Linking}). Este método identifica las
entidades mencionadas en la pregunta del usuario y las asocia correctamente con
las tablas y columnas de la base de datos, reduciendo las confusiones
\citep{qin2024gensql}. Estudios recientes sugieren el uso de redes y análisis
de atención para mejorar esta conexión. Esto ayuda a que el modelo comprenda
cómo se relacionan los datos entre sí antes de empezar a escribir el código SQL
\citep{gladkykh2025datrics}.

Cuando las consultas se vuelven más complejas, los sistemas modernos recurren a
estrategias de razonamiento paso a paso. En lugar de generar todo el código de
una sola vez, los modelos dividen la pregunta en problemas más pequeños y
validan cada paso antes de dar la respuesta final \citep{shuai2025deepvis}.
Esta estrategia, sumada al uso de modelos entrenados específicamente para
programar, aumenta notablemente la precisión de los resultados en entornos
corporativos reales que manejan grandes volúmenes de datos
\citep{lei2024spider}.

\subsection{Generación Aumentada por Recuperación (RAG)}

Para superar los límites de memoria de los modelos de lenguaje y evitar que
inventen información, se utiliza una arquitectura llamada Generación Aumentada
por Recuperación (RAG). Esta técnica permite que la Inteligencia Artificial
consulte documentos externos antes de generar una respuesta, reduciendo su
dependencia exclusiva de los datos con los que fue entrenada originalmente. En
entornos corporativos, esto es fundamental para garantizar que el modelo
responda basándose únicamente en la información validada de la propia empresa
\citep{arguello2025rag}.

El funcionamiento interno de RAG se divide en dos fases: recuperación y
generación. Primero, la información de la empresa se transforma en
representaciones matemáticas llamadas vectores y se guarda en una base de datos
especial. Cuando un usuario hace una pregunta, el sistema busca matemáticamente
cuáles son los fragmentos de información más similares. Luego, esta información
recuperada se le entrega al modelo de lenguaje para que tenga el contexto exacto
antes de formular su respuesta \citep{gao2024retrieval}.

En el análisis de datos, RAG se emplea para buscar dinámicamente descripciones
de tablas o ejemplos de consultas que ya han sido resueltas por expertos. Al
incluir estos ejemplos exitosos junto con la pregunta del usuario, el modelo
mejora significativamente su capacidad para entender el problema y redactar el
código correcto. Este enfoque disminuye los errores y facilita la resolución
automática de preguntas analíticas con un alto grado de precisión
\citep{LeveragingNLPAndRAGModelsForAutomatedBusinessIntelligenceQueryResolution}.

\subsection{Sistemas Multi-Agente}

El avance de la Inteligencia Artificial ha permitido pasar de modelos que
simplemente responden preguntas a agentes autónomos. Un agente se define por su
capacidad de mantener un objetivo, recordar conversaciones pasadas, planificar
tareas paso a paso y utilizar herramientas externas, como ejecutar código o
consultar bases de datos. Esta capacidad de razonar y actuar de forma
independiente permite resolver problemas mucho más complejos que la simple
generación de texto \citep{maddila2024agentes}.

En lugar de depender de un solo agente para hacer todo el trabajo, las
arquitecturas modernas utilizan sistemas multi-agente. En este esquema, el
problema se divide entre varios especialistas virtuales que colaboran entre sí.
Por ejemplo, un agente puede encargarse exclusivamente de escribir el código
SQL, mientras que otro actúa como revisor para detectar errores sintácticos
antes de ejecutar la consulta. Esta división del trabajo imita el flujo de
desarrollo de un equipo humano y aumenta notablemente la calidad del resultado
final \citep{wang2024macsql}.

En el contexto de la Inteligencia de Negocios, la orquestación de múltiples
agentes representa un pilar para la automatización segura. Al distribuir la
responsabilidad, los sistemas pueden enfrentar de manera colaborativa los
desafíos de interactuar con bases de datos empresariales. Esta supervisión
cruzada reduce los errores del modelo y asegura que las decisiones analíticas
se basen en datos precisos y correctamente procesados
\citep{gartner2025predicciones}.

\subsection{Generación automática de visualizaciones}

En la Inteligencia de Negocios, la representación gráfica de los datos es el
paso final indispensable tras realizar una consulta. La investigación actual
busca que los modelos de lenguaje no solo devuelvan tablas con números, sino
que traduzcan las peticiones del usuario directamente en visualizaciones de
datos. El objetivo es que personas sin conocimientos técnicos puedan crear y
consumir gráficos interactivos simplemente describiendo lo que necesitan ver
\citep{perez2026ia}.

Para lograr que la Inteligencia Artificial diseñe gráficos precisos, se utilizan
metodologías que obligan al modelo a razonar su proceso creativo. Antes de
escribir el código de visualización, el sistema planifica lógicamente qué tipo
de gráfico es el más adecuado para los datos obtenidos, qué variables irán en
cada eje y qué paleta de colores utilizar. Este razonamiento paso a paso mejora
la relación entre la pregunta original y el diseño final
\citep{shuai2025deepvis}.

A pesar de estos avances, los sistemas actuales suelen utilizar herramientas de
código ligeras para generar los tableros, lo que facilita el desarrollo pero
introduce nuevas limitaciones
\citep{shi2026nl2dashboardlightweightcontrollableframework}. Problemas como
escalas incorrectas en los ejes o la falta de opciones avanzadas de
interactividad son comunes. Por lo tanto, mejorar la fidelidad gráfica y lograr
que estos tableros compitan con las herramientas tradicionales sigue siendo un
área que requiere mayor investigación \citep{perez2026ia}.

\subsection{Evaluación y modelos como jueces (LLM-as-a-Judge)}

Evaluar la calidad de un sistema de Inteligencia de Negocios Generativa (GenBI)
es una tarea sumamente compleja, ya que el proceso abarca desde la extracción de
datos hasta el diseño visual. Tradicionalmente, la evaluación se limitaba a la
etapa de generación de consultas (Text-to-SQL), midiendo si el código escrito
por el modelo era idéntico a una respuesta esperada. Como esta métrica resultó
demasiado rígida, la industria migró hacia medir la precisión de ejecución
(\textit{Execution Accuracy}), verificando únicamente si los datos devueltos
son correctos \citep{lei2024spider}. No obstante, cuando el sistema también
genera tableros y gráficos, estas métricas matemáticas se vuelven insuficientes.

Dado que es imposible revisar manualmente miles de resultados que combinan
código SQL, estructuras de datos y representaciones visuales, ha cobrado fuerza
la estrategia de utilizar modelos de lenguaje como jueces automatizados
(\textit{LLM-as-a-Judge}). En este enfoque, se emplea un modelo avanzado para
leer la pregunta original del usuario y evaluar semánticamente si todo el flujo
tiene sentido. El modelo actúa como un auditor humano, verificando que el
razonamiento detrás de la respuesta sea lógicamente válido y útil para el
negocio \citep{Jaisinghani2024}.

Esta metodología de evaluación resulta indispensable para auditar ciclos
analíticos de principio a fin. Al utilizar un modelo como juez, el sistema
puede calificar de manera integral tanto la exactitud de los datos extraídos
como la coherencia de la explicación textual y la claridad del gráfico generado.
Contar con este nivel de evaluación flexible permite realizar análisis
comparativos sólidos entre diferentes arquitecturas, garantizando que el tablero
final cumpla con los estándares de calidad corporativos \citep{Jaisinghani2024}.